\chapter{Introduction}
\label{ch:intro}

% \section*{Background and Motivation}

The standard framework of statistical learning deals with extracting underlying patterns from identically structured, stationary empirical observations~\cite{vapnik1998statistical,shalev-shwartz_ben-david_2014}. Under this baseline, past observations guide future actions with rigorous generalization guarantees. However, there exists an important class of real-world decision-making problems that depart from these core assumptions. This dissertation is motivated by the need to systematically study and extend statistical learning theory to two such distinct operational departures.

This dissertation investigates two key extensions of the classical baseline. The first extension considers strategic environments featuring an intelligent agent who actively influences the data-generating law, for instance, by manipulating feature values or selectively withholding specific feature components~\cite{hardt2016strategic,grossman1981informational,milgrom1981good}. Consequently, statistical patterns learned from historical observations may no longer remain valid once the agent adapts to the deployed decision rule. Any effective data-driven decision-making procedure must therefore explicitly account for the incentive structure of the strategic agent in order to remain effective.

The second extension considers settings where observations depart structurally from a uniform feature baseline: predictive features and labels are partitioned across distinct training sets and never jointly observed during training~\cite{little2019statistical,vapnik2009new}. This represents a structural information constraint rather than an agent incentive problem. In the classical setting, collecting more samples simply increases the quantity of information. However, when observation structures vary, different samples contribute distinct qualities and quantities of information. Consequently, the total information available for learning depends fundamentally not only on the number of observations, but also on which variables are jointly observed.

The central objective unifying this dissertation is to characterize achievable performance and the fundamental limitations of learning under these two distinct departures. For strategic learning, we address core questions such as: \emph{What is the optimal decision rule when every parameter of the environment is known (referred to as the decision problem in this dissertation)? How can this optimal structure be learned from finite data? How does the resulting performance compare with the non-strategic baseline?} For learning with partitioned information, we characterize the fundamental learning limits imposed by which variables are jointly observed, develop an imputation-based empirical risk minimization framework for such partitioned settings, and evaluate how the choice of imputation governs learning performance relative to the benchmark where complete joint observations are available.

To address these challenges, we adopt a unified two-stage methodology throughout the dissertation. We first analyze the underlying decision problem assuming complete knowledge of the environment, which unveils key structural properties of the optimal decision rule. These structural insights then directly guide the design of the learning algorithm. Rather than searching over a large, unstructured hypothesis space, the learning procedure searches over a structured space suggested by the analysis of the decision problem, yielding more sample-efficient and accurate learning. For strategic learning, the natural mathematical framework is game theory, specifically sender--receiver games where the decision maker acts as the leader committing to a rule, commonly known as screening games, evaluated under the Stackelberg equilibrium concept~\cite{bruckner2011stackelberg,hardt2016strategic}. Incorporating this game-theoretic structure explicitly accounts for the strategic behavior of agents during learning. The theoretical results in this dissertation either characterize structural properties of the optimal decision rule or reformulate the underlying decision problem into a structured optimization problem that is easier to analyze. For learning with partitioned observations, the central challenge arises from statistical dependencies that are never jointly observed during training. We address this through a family of imputation-based learning algorithms, where the choice of imputation directly determines the achievable learning performance.

The mathematical models studied in this dissertation are intentionally formulated with structural simplicity to maintain analytical tractability and physical interpretability. Nevertheless, the theoretical insights and results derived from these models serve as valuable guiding principles for more complex settings where exact analytical solutions may be unavailable. A recurring theme that emerges across strategic learning problems is that strategic behavior need not always be viewed purely as an obstacle to be eliminated. When modeled explicitly, strategic actions can reveal additional information and, in certain settings, be leveraged to improve overall decision making.

\section*{Organization of the Dissertation}
\label{sec:intro_contrib}

The chapters of this dissertation are organized as follows. Each chapter is self-contained and presents independent research contributions that can be read without prerequisite dependencies on the others.

Chapter~\ref{ch:strat_class} studies the strategic classification problem where agents have non-uniform preferences over labels. Existing works typically assume that all agents share a uniform preference for a single label. Relaxing this assumption fundamentally changes the structure of the strategic classification problem. The chapter investigates how strategic decision rules should differ from their non-strategic counterparts and whether deliberately making classification errors near the class boundary can reduce the overall strategic classification error. It then studies the associated learning problem and shows that explicitly incorporating the underlying game-theoretic structure can improve learnability in certain situations. The work presented in this chapter is based on our paper published in \emph{Autonomous Agents and Multi-Agent Systems} (JAAMAS, 2025)~\cite{singh2025strategic}, with preliminary portions appearing in NCC~2024~\cite{singh2024optimal} and ICC~2022~\cite{singh22strategic}.

Chapter~\ref{ch:side_info} studies learning from partitioned datasets, where the target variable and the side information are never jointly observed during training but are both available at test time. Unlike the classical missing-feature setting, in which a limited number of joint observations still exist, the partitioned information structure prevents direct learning of the relationships between the target variable, the features, and the side information. The chapter characterizes the fundamental learning limits imposed by this information structure, examines a family of imputation-based learning algorithms, and investigates how the choice of imputation governs the achievable learning performance. It further shows that, under suitable conditions, the optimal bound-minimizing regularization strategy can be restricted to a small discrete set of candidates. The results in this chapter are based on our manuscript submitted to \emph{Information and Inference: A Journal of the IMA} (Oxford University Press).

Chapter~\ref{ch:strat_rev} studies a regression problem in which a strategic feature provider selectively reveals only a subset of its verifiable features to the learner. Unlike strategic manipulation, where feature values are altered, the strategic action here lies in the choice of which features to disclose. The chapter establishes that, in certain cases, the information conveyed by the disclosure decision, together with the revealed features, can improve prediction performance relative to the non-strategic baseline. It then characterizes the optimal disclosure policy and studies the corresponding offline and online learning problems arising from this strategic interaction. The work in this chapter is based on our paper accepted at the \emph{IEEE Conference on Decision and Control} (CDC, 2026)~\cite{singh2026linear} and our manuscript submitted for publication~\cite{singh2026revelation_journal}.

Chapter~\ref{ch:ipo} deals with the regulation of Initial Public Offerings (IPOs)~\cite{ritter2003investment}, where the firm can strategically window dress its presentation to influence investor decisions, while the regulator determines the extent to which such manipulation should be permitted. The chapter establishes that, in certain cases, allowing a limited degree of strategic window dressing can improve market outcomes by correcting investors' bias against fundamentally good firms. It then studies the corresponding learning problem, where the regulator learns effective regulatory policies when the underlying market environment is unknown. The results in this chapter are based on our manuscript submitted for publication~\cite{singh2026ipo_journal}.

\clearpage
\section*{Publications}
\label{sec:publications}

\subsubsection*{Journal Articles:}
\begin{itemize}[leftmargin=*,noitemsep,topsep=2pt]
\item[\(1.\)] M. K. Singh and A. A. Kulkarni. ``Strategic Classification for Non-uniform Preferences using Penalty Labels and Randomisation,'' \emph{Autonomous Agents and Multi-Agent Systems (JAAMAS)}, 2025.
\item[\(2.\)] M. K. Singh, S. S. Kowshik, and A. A. Kulkarni. ``Learning with Side Information from Partitioned Data,'' \emph{Information and Inference: A Journal of the IMA} (Oxford University Press), (\emph{Submitted}).
\item[\(3.\)] M. K. Singh, A. A. Kulkarni, and K. Mahajan. ``Learning under Strategic Feature Revelation,'' (\emph{Submitted for publication, 2026}).
\item[\(4.\)] M. K. Singh and A. A. Kulkarni. ``Strategic IPO Regulation under Complete and Incomplete Information,'' (\emph{Submitted for publication, 2026}).
\end{itemize}

\subsubsection*{Conference Articles:}
\begin{itemize}[leftmargin=*,noitemsep,topsep=2pt]
\item[\(1.\)] M. K. Singh, A. A. Kulkarni, and K. Mahajan. ``Learning under Strategic Feature Revelation,'' \emph{IEEE Conference on Decision and Control (CDC)}, 2026 (\emph{Accepted}).
\item[\(2.\)] M. K. Singh and A. A. Kulkarni. ``Optimal Stochastic Decision Rule for Strategic Classification,'' \emph{National Conference on Communications (NCC)}, 2024.
\item[\(3.\)] M. K. Singh and A. A. Kulkarni. ``Strategic Multiclass Classification with Non-uniform Preferences and its Relation to Incentive Compatibility,'' \emph{Indian Control Conference (ICC)}, 2022.
\end{itemize}
